% Pilot Fabric Live Companion and Private Continuity
% Overleaf dependencies: \usepackage{booktabs,array,hyperref}
% This file is a section intended to be included in the main AION/Pilot document.

\section{Pilot Fabric: Live Companion, Private Continuity, and the Television-to-Agent Wedge}
\label{sec:pilot-live-companion-private-continuity}

\subsection{Purpose and Product Position}

Pilot Fabric turns an existing television into the room-scale surface of a
governed household intelligence without requiring a replacement television or
new proprietary hardware.  The television is the adoption wedge: users first
encounter Pilot through immediately understandable capabilities such as media
control, contextual explanation, translation, sport interpretation, games,
learning, and saving something seen on screen.  The same installed intelligence
can subsequently continue privately on a phone and, after explicit approval,
be extended into personal tasks, services, vehicles, contacts, and other homes.

The central product rule is therefore:

\begin{quote}
The shared television may perceive, explain, prepare, and propose.  Identity,
private memory, recipients, credentials, purchases, bookings, messages, and
other consequential actions must cross a private, persona-bound confirmation
boundary before execution.
\end{quote}

Pilot Fabric releases 0.22--0.26 establish this boundary across sharing, scene
explanation, translation, sport interpretation, and programme-to-private
saving.  They build on the consumer reliability foundation introduced in
release 0.21: persistent controller addresses, recoverable TV connectivity,
automatic process recovery, and an isolated runtime that does not overwrite
COMDEX or AION's existing intelligence files.

\subsection{System Architecture}

The implementation represents one governed intelligence across four different
surfaces:

\begin{enumerate}
    \item \textbf{Mother brain.}  A capable computer hosts the cognitive runtime,
    local speech recognition, local models, policy enforcement, memory, provider
    adapters, audit ledger, and orchestration loop.

    \item \textbf{Television gateway.}  A restricted television such as an LG
    webOS device is represented by a signed gateway.  The gateway exposes only
    discovered and permissioned capabilities, for example observing application
    state, reading volume, launching an allowlisted surface, or sending bounded
    navigation controls.  The TV never receives mother-brain secrets.

    \item \textbf{Pilot Canvas.}  A token-protected television web surface
    displays Pilot experiences, evidence, learning activities, God View, and
    pairing information.  Contextual answers normally preserve the programme
    instead of replacing it with a full-screen result.

    \item \textbf{Private phone companion.}  The phone provides the private
    control and approval boundary.  It displays detailed evidence, identity,
    recipients, proposed service scopes, saved items, and explicit approval
    controls.  Requests return to the mother brain; TV pairing keys, COMDEX
    credentials, and provider secrets are not projected into the phone page.
\end{enumerate}

The logical data path is:

\begin{center}
\small
Shared voice or phone request $\rightarrow$ signed capability check
$\rightarrow$ bounded live context $\rightarrow$ deterministic/local
intelligence $\rightarrow$ confidence and evidence validation
$\rightarrow$ spoken TV response $+$ private phone record
$\rightarrow$ optional private approval $\rightarrow$ verified receipt.
\end{center}

Every stage is policy constrained.  A model may propose structured content, but
it does not directly operate a device or external account.  Device operations
are issued by deterministic adapters after capability validation.  External
actions require a persona-bound adapter, an exact private preview, approval, and
a verifiable outcome receipt.

\subsection{Bounded Live Context and Perception}

An explicit Pilot request opens a bounded live-context window.  The mother
combines the following permitted evidence:

\begin{itemize}
    \item current webOS application and audio state;
    \item a short, in-memory window of locally transcribed dialogue;
    \item owner-initiated phone-camera OCR and classification;
    \item structured subtitle, scoreboard, product, location, and scene findings;
    \item owner corrections associated with the current observation;
    \item bounded provider metadata, when an official or permitted identifier is
    available; and
    \item the current task and device-belief state.
\end{itemize}

Source images are processed locally and deleted immediately.  Raw microphone
audio remains ephemeral and is cleared when voice is put to sleep.  The durable
record contains structured observations, source types, confidence, timestamps,
and cryptographic hashes rather than retained programme pixels or raw audio.
Programme dialogue is evidence, never an instruction to Pilot.

A simplified live-context record has the following shape:

\begin{verbatim}
{
  "schema_version": "pilot.live.context.v1",
  "device_state": {"app": "...", "volume": 18},
  "recent_transcripts": ["bounded local dialogue"],
  "screen_observation_id": "observation_...",
  "evidence_types": ["device_state", "dialogue", "owner_capture"],
  "context_hash": "sha256...",
  "raw_audio_retained": false,
  "source_pixels_retained": false
}
\end{verbatim}

The context hash binds downstream explanations, translations, Moments, sport
interpretations, and private-save cards to the evidence used to create them.

\subsection{Pilot Moments and Safe Sharing: Release 0.22}

Pilot Moments provide a rights-aware method for sharing the meaning or evidence
of something on television without copying protected programme video or audio.
For example, after completing a live fact-check, the command
``Pilot, send this fact-check'' prepares a signed card containing the exact
claim, verdict, confidence, source links, and evidence hash.

Preparation does not send anything.  The recipient is selected only on the
private phone and is not spoken or displayed on the shared television.  Binding
a recipient produces an immutable preview hash.  The user must then press a
separate \emph{Send this exact Moment} control.  Any change to the card or
recipient invalidates that preview and requires confirmation again.

The Moment lifecycle includes:

\begin{itemize}
    \item signed evidence and issuer identity;
    \item private-phone-only recipient selection;
    \item an exact preview followed by a separate Send decision;
    \item adapter-gated delivery and a verified delivery receipt;
    \item an honest \texttt{approved\_pending\_adapter} state when no authorized
    account is connected;
    \item 24-hour expiry, replay protection, bounded delivery attempts, and rate
    limiting;
    \item revoke, private-data deletion, and abuse reporting; and
    \item privacy-preserving audits that store recipient hashes rather than raw
    addresses.
\end{itemize}

No successful-delivery claim is permitted without a receipt from an authorized
adapter.  A provider-native timestamp or clip mechanism can later replace the
context-card fallback where the provider officially permits it.

\subsection{Spoiler-Aware Scene Explanation: Release 0.23}

Pilot distinguishes scene explanation from general web research.  Commands such
as ``what just happened?'', ``what did they just say?'', and ``explain this scene
without spoilers'' construct an observed-evidence envelope from recent local
dialogue, captured subtitles, screen observations, and owner corrections.

Public plot summaries, endings, future events, and unseen character knowledge
are deliberately excluded.  A local Gemma model may interpret the envelope,
but its response passes through an independent future-event leakage guard.  If
the model is unavailable or its result is unsafe, Pilot uses a deterministic
observed-evidence fallback.  When the evidence is insufficient, Pilot refuses
to guess and asks for a new observation or a request made immediately after the
relevant dialogue.

The television keeps playing.  Pilot speaks a short explanation while the
phone displays the observed facts, confidence, uncertainty, evidence types, and
record hash.  This capability requires no OpenAI API.

\subsection{Local Live Translation: Release 0.24}

Live Translation selects only the bounded recent line or an explicitly
requested, owner-captured subtitle.  It supports English, Spanish, French,
German, Italian, and Portuguese.  The routing order is local Gemma first,
followed by a deterministic offline phrasebook for supported common expressions;
it does not silently invoke a paid model.

The model receives a translation-only schema that prohibits explanation,
continuation, added context, and instruction execution.  A validator outside
the model requires protected tokens---including numbers, times, and URLs---to
survive unchanged.  If no reliable local result exists, Pilot states the
limitation instead of inventing a translation.  A language-appropriate installed
system voice speaks the short result while playback continues.

One field test translated ``Nos vemos en la estaci\'on a las 8:30'' as ``See you
at the station at 8:30.''  The time survived independent validation exactly,
with no internet or paid-AI dependency.

\subsection{Evidence-Bounded Live Sports: Release 0.25}

Live Sports structures teams, scores, and match clocks from the most recent
owner-captured scoreboard.  It can state which side leads, but labels the score
as a captured snapshot that may become stale.  Commentary and subtitle cues are
treated as observations; rule descriptions are separately labelled inference.

A deterministic local rule library covers football offside, handball,
penalties, fouls, cards, and VAR; basketball travelling and contact; and tennis
tie-breaks.  Pilot never identifies an unseen player and never claims a referee
was correct or incorrect merely because a rule matches a named incident.  An
ambiguous question such as ``explain that referee decision'' is refused when no
incident cue was captured.  The match remains visible while the detailed score,
rule, evidence, confidence, and limitation appear privately on the phone.

This release establishes a critical distinction for future live feeds:

\begin{quote}
Screen and commentary evidence describe what Pilot observed.  Rule knowledge
explains what a named rule means.  Only an authenticated live-data provider may
establish current official statistics or event outcomes.
\end{quote}

\subsection{Programme-to-Private Saving: Release 0.26}

Commands such as ``save this product'', ``remember this recipe'', ``save this
destination'', ``save this song'', and ``add this idea to my phone'' create a
short-lived evidence card.  Candidate content is assembled from bounded dialogue,
the latest screen observation, product or location findings, and provider
metadata.  Each item contains a category, title, detail, evidence lines,
confidence, creation and expiry times, and evidence and item hashes.

The shared TV does not know which household persona owns the request.  A pending
card therefore has no \texttt{persona\_id} and remains in
\texttt{awaiting\_private\_claim} for at most 30 minutes.  The phone displays the
card under \emph{Save from TV}.  Only pressing \emph{Save to my Pilot} binds it
to the active private identity.

\begin{verbatim}
{
  "schema_version": "pilot.private-save.v1",
  "category": "product",
  "status": "awaiting_private_claim",
  "persona_id": null,
  "confidence": 0.91,
  "evidence_hash": "sha256...",
  "privacy": {
    "shared_tv_selected_persona": false,
    "private_claim_required": true,
    "purchase_executed": false,
    "message_sent": false,
    "raw_audio_retained": false,
    "source_pixels_retained": false
  }
}
\end{verbatim}

Saved collections are filtered by persona.  Another identity cannot view the
identity-bound pending record, inspect the saved collection, or delete its
items.  Saving is deliberately a memory action only: it does not purchase,
book, message, create a calendar event, or choose an external account.  Those
operations require a later, separately scoped research and service proposal.

\subsection{Local-First Intelligence Routing}

These capabilities follow the AION Native routing policy:

\begin{enumerate}
    \item deterministic local capability;
    \item existing AION knowledge, memory, and governed cognitive runtime;
    \item an installed local model;
    \item bounded free or public evidence retrieval;
    \item optional user-supplied Gemini synthesis or grounding;
    \item optional premium providers, including OpenAI; and
    \item an honest limitation.
\end{enumerate}

Cloud providers improve speed or quality but are not allowed to remove Pilot's
core usefulness when absent.  Results identify their route, for example
\emph{AION Local}, \emph{AION plus live public evidence}, or \emph{Pilot Boost}.
Retrieval, citation validation, contradiction detection, action permission, and
receipt verification remain outside the generative model.

\subsection{Security, Identity, and Audit Properties}

The releases preserve the following invariants:

\begin{itemize}
    \item device actions require enrollment and signed capability capsules;
    \item nodes and gateways possess independent Ed25519 identities;
    \item capability scope, risk, expiry, sequence, and replay state are checked
    before execution;
    \item shared speech cannot select a private persona, recipient, card, calendar,
    service account, or payment reference;
    \item provider passwords, tokens, private keys, full card numbers, CVVs, PINs,
    and raw credentials are rejected from ordinary memory;
    \item protected programme audio and video are not copied into Moments;
    \item source phone-camera frames are deleted after local interpretation;
    \item audit records store bounded metadata and hashes rather than raw private
    content; and
    \item the append-only audit chain is verified on every release field test.
\end{itemize}

The universal-node model permits one installed node to trust multiple mother
brains through separate stored handshakes and short-lived capability leases.
Each mother and persona retains isolated preferences, memory, accounts, and
approvals.  A second resident can reuse a compatible installed node without
silently inheriting the first resident's private state.

\subsection{Functional Command Surface}

\begin{table}[htbp]
\centering
\small
\begin{tabular}{p{0.24\linewidth} p{0.68\linewidth}}
\toprule
Capability & Example requests \\
\midrule
Safe evidence sharing & ``Pilot, send this fact-check''; ``Pilot, share that Moment'' \\
Scene explanation & ``Pilot, what just happened?''; ``Pilot, explain without spoilers'' \\
Local translation & ``Pilot, translate that into Spanish''; ``translate that subtitle to English'' \\
Live sports & ``Pilot, what's the score?''; ``who is winning?''; ``why was that offside?'' \\
Private saving & ``Pilot, save this product''; ``remember this recipe''; ``add this idea to my phone'' \\
\bottomrule
\end{tabular}
\caption{Live Companion and private-continuity command examples.}
\label{tab:pilot-live-command-surface}
\end{table}

\subsection{Verification and Release Evidence}

\begin{table}[htbp]
\centering
\small
\begin{tabular}{p{0.11\linewidth} p{0.44\linewidth} p{0.14\linewidth} p{0.16\linewidth}}
\toprule
Release & Principal capability & Automated checks & Field state \\
\midrule
0.22 & Private confirmed Moment sharing & 180 & LG connected \\
0.23 & Spoiler-aware scene explanation & 182 & LG connected \\
0.24 & Local validated translation & 186 & LG connected \\
0.25 & Evidence-bounded sport interpretation & 190 & LG connected \\
0.26 & Persona-bound programme saving & 195 & LG browser observed; audit valid \\
\bottomrule
\end{tabular}
\caption{Incremental Fabric verification through release 0.26.}
\label{tab:pilot-fabric-release-evidence}
\end{table}

Release 0.26 was packaged as
\texttt{aion-fabric-preview-0.26.0.zip} with SHA-256 digest
\texttt{680d051c1037049e\allowbreak
71069b1b2645b6c8\allowbreak
c2ecf2289fc267e8\allowbreak
7170791ad22f69ab}.
The live field check observed the paired LG webOS browser at volume 18, verified
the audit chain, and confirmed that the persistent private phone connection
survived the upgrade.  Existing COMDEX intelligence remained untouched.

\subsection{Remaining Engineering Roadmap}

The remaining build order is intentionally strict.

\subsubsection{Pilot Fabric 0.27: Live Event Intelligence}

The next release should connect authenticated sports and live-event feeds for
scores, fixtures, player statistics, event timelines, elections, concerts, and
awards.  Provider data must be fused with current screen perception while
preserving the distinction between an official live fact, a captured screen
observation, and Pilot's interpretation.  Intended requests include ``who
scored?'', ``show the statistics'', and ``notify me if it becomes close''.

\subsubsection{Saved-Item Follow-Through}

Privately saved products, recipes, destinations, songs, and ideas should become
launch points for further work.  Pilot may research, compare, create a shortlist,
or prepare a task, trip, playlist, shopping proposal, or learning activity.  A
saved item must not itself authorize a purchase, booking, message, or calendar
change; each consequential transition requires a new exact scope and private
approval.

\subsubsection{Verified Entertainment and Navigation}

Pilot should add provider-verified deep links, reliable title search, profile
selection, playback, and ``continue what I was watching''.  Subscription state
must be obtained through authorized account adapters rather than inferred from
search results.  Every third-party operation should implement
observe--act--observe--verify, automatic recovery, alternative paths, duplicate
suppression, and per-device memory of successful and failed navigation.

\subsubsection{Production Perception and Consumer Quality}

The phone observer requires production TLS or a signed native shell, reliable
reconnection after sleep, multi-frame temporal understanding, and authenticated
service metadata.  Product completion additionally requires native LG webOS and
equivalent Samsung and Android/Google TV surfaces, multi-room TV routing, stable
focus and Back/Home behaviour, better background-dialogue rejection, natural
selectable voices, interruption and cancellation, automatic updates, crash
recovery, diagnostics, accessibility review, and child-safety review.

Games should gain controller detection, verified launch, saved shortcuts, and
``continue last game''.  Learning should add a broader curriculum, licensed
visuals, pronunciation assessment, parent-visible progress, difficulty bands,
and additional subjects.

\subsubsection{Personal Pilot and the Natural Network}

After the television experience is dependable, Pilot should introduce native
signed-phone identity, biometric approvals, memory inspection/correction/export/
deletion, and cross-device task continuation.  The Life Concierge then connects
to-dos, reminders, calendars, email, messaging, maps, shopping, booking, music,
car continuation, and carefully configured emergency support through the
sequence research--prepare--private review--approve--execute once--verify--receipt.

The final network stage adds privately confirmed contacts, signed Pilot-to-Pilot
requests, family coordination, shared lists, useful web Moments for non-users,
and optional invitations.  No request may silently assign a task, disclose a
location, send a message, make a purchase, or create a booking.  Growth must
result from a useful artifact and explicit user choice rather than forced
invitations or dark patterns.

\subsection{Conclusion}

Releases 0.22--0.26 prove more than a collection of television commands.  They
establish a reusable architecture in which room-scale shared intelligence can
observe and explain current content, preserve programme playback, continue
privately on a phone, bind data to the correct identity, and stop before
consequential execution.  This creates the bridge from a compelling AI-TV
experience to a general personal agent and, ultimately, a trusted agent network.
